\documentclass[11pt]{article}
\usepackage[margin=1in]{geometry}
\usepackage{amsmath,amssymb,hyperref}
\newcommand{\F}{\mathbb F}
\newcommand{\Res}{\operatorname{Res}}
\newcommand{\E}{\mathbb E}
\newcommand{\wtO}{\widetilde O}
\title{Direct evaluation and differentiation of the scalar target norm}
\author{Project research note}
\date{18 November 2026}
\begin{document}
\maketitle

One homogeneous univariate resultant evaluates the normalized scalar
target norm without expanding the product of all the slots.
On an isolated tuple the same split gives every requested first
derivative. A separate two-slot calculation gives the full additive
gradient even at a nonzero norm. At the tuned widths the work needed
by the relation finder drops from a dense root algebra of dimension
$2^{(1-o(1))b}$ to an explicit resultant with work
$2^{(1/2+o(1))b}$. The existing conditional end-to-end exponent
is smaller asymptotically. We also test whether the arithmetic indices
of the actual dyadic masks supply reusable elliptic translations.

\section{The exact scalar split}
Use $E:y^2=x^3+ax+c$ over $\F_p$, $p>3$, and $t=x(T)$ with
$2T\ne O$. Initially each slot contains both rational lifts of
each abscissa in $U_i$, with no fixed points of negation. Let
$n_i=|U_i|$. An empty slot gives norm one and zero derivatives.
For one abscissa tuple $\mathbf u$ write $S(\mathbf u;Z)$ for
the normalized homogeneous sign-orbit form constructed in
\texttt{../joint\_parity/parity.tex}. Its formal degree in $Z$ is
$2^{m-1}$. Specifically, start with $Z-u_1$ and append each
abscissa $v$ using the homogeneous resultant with
\[
 K(W,v,Z)=(W-v)^2 Z^2-2((W+v)(Wv+a)+2c)Z+(Wv-a)^2-4c(W+v).
\]
Homogenize separately in each argument. The convention for a finite
input is $(u:1)$, and for infinity it is $(1:0)$. The scalar and
additive derivatives requested here refer to the particular product
\[
 N(z,\mathbf s)=\prod_{\mathbf u}
  \big(S(\mathbf u;z)+s_0+\textstyle\sum_i s_i u_i\big)
\]
at $(z,\mathbf s)=(t,\mathbf0)$. The normalization matters for
these absolute derivatives.

For $m\ge2$ split the slots into nonempty blocks $L,R$, initially
a prefix and its complement. For each left tuple build its forward
form, with output coordinate $W$. For each right tuple start with
$t$ and append its slots in reverse order, again with output $W$.
Multiply within each block to obtain
\[
 A(W)=\prod_{\mathbf l} S(\mathbf l;W),\qquad
 B(W)=\prod_{\mathbf r} S(t,\operatorname{rev}\mathbf r;W).
\]
Writing $D_I=\prod_{i\in I} n_i$, their formal degrees are
\[
 f=2^{|L|-1}D_L,\qquad h=2^{|R|}D_R .
\]
Then $N(t,\mathbf0)=\Res_{f,h}(A,B)$ exactly.

To see the scalar identity, factor the forward form over an algebraic
closure into homogeneous linear factors, retaining its leading
scalar. One more chord step sends each root $U$ to the quadratic
$K(U,v,V)$ and squares the leading scalar. If $G(V)$ is the form
coming from the remaining backwards steps, the contribution of this
quadratic is $\Res_V(K(U,v,V),G(V))$. Multiplicativity of the
homogeneous resultant gives the same expression as eliminating $U$
first. In the reverse recurrence one interchanges these two resultant
arguments. A quadratic times any integer degree is even, so the
interchange has positive sign. The last backwards step is
$K(U,v,t)=K(t,v,U)$. This proves the identity for a pair of tuples,
including all scalars. It is a polynomial identity and includes
infinite roots. Multiplicativity over all tuples gives the stated
single resultant. The two block products can be built by subproduct
trees. Elliptic biquadratic product evaluations by resultants also
appear in \href{https://velusqrt.isogeny.org/velusqrt-20200616.pdf}
{Bernstein, De Feo, Leroux and Smith, Sections 2 and 4}.

For later use the normalized form is symmetric in its $m+1$
arguments when $m\ge2$. One verification uses the irreducible
hypersurface which is the image of $E^m$ under the coordinate maps
and the abscissa of the sum. Its set of points is invariant under
permutations, because the signed equation is a zero sum. The
projective biquadratic has no identically zero specialization on
a pair of inputs, including infinity and rational two-torsion.
Thus there are no extra components. Generic sign orbits have
distinct final abscissas, so the hypersurface occurs with
multiplicity one. The same multidegree $2^{m-1}$ is obtained in
each argument by the recurrence. The scalar under permutation is
one by evaluating with one finite argument and all others at
infinity. The normalized multihomogeneous form takes value one at
each of these specializations by the same recurrence. This also
permits any choice of split. See the original
\href{https://eprint.iacr.org/2004/031}{summation-polynomial construction}.

\section{Differentiate the resultant at a common root}
First hold $A$ fixed and vary $B$ in direction $V$, preserving
formal degree $h$. Let the ordinary affine degree of $A$ be $n$,
its leading coefficient be $a_0$, $e=f-n$ and $F=A/a_0$ in
the affine chart. Put $b_h=[W^h]B$. The homogeneous formula is
\[
 \Res_{f,h}(A,B)=(-1)^{eh}a_0^h b_h^e\,R,\qquad
 R=\Res(F,B)=\prod_{F(u)=0}B(u),
\]
with multiplicity in the product and $R=1$ for $n=0$.
The sign comes from $\Res(W_1,B)=(-1)^h B(1,0)$.
Differentiate $b_h^e$ by the ordinary power rule, including when
$b_h=0$. The finite product has the following derivative.

\begin{itemize}
\item If $F=1$ its derivative is zero. If $\gcd(F,B)=1$ and
$n>0$, compute $J=B^{-1}F'\bmod F$. Then
$\delta R=R[W^{n-1}](VJ\bmod F)$.
Indeed the coefficient is the trace of multiplication by $V/B$
in $\F_p[W]/F$. This follows from the coefficient of $W^{-1}$
in $(V/B\bmod F)\,F'/F$ at infinity. Each root contributes its
multiplicity, also in positive characteristic.
\item If the monic gcd is $W-u$ and $F=(W-u)F_1$ with
$F_1(u)\ne0$, then $\delta R=\Res(F_1,B)V(u)$. If more than
one factor in the root product vanishes, $\delta R=0$. This
includes a repeated root of $F$ even when the gcd has degree one.
\end{itemize}

For a simultaneous change of $A$ and $B$ add the contribution with
the roles swapped and the sign $(-1)^{fh}$. Consequently products,
extended gcds and remainders on two univariate polynomials suffice
even when the resultant vanishes. The verifier uses the corresponding
\href{https://flintlib.org/doc/algorithms/polynomials/gcd.html}{FLINT polynomial operations}.

Differentiate the half forms by the product rule. In addition to the
$z$ derivative, compute two public polynomial deformations for slot
$i$, moving every one of its abscissas by $u\mapsto u+\epsilon$ or
$u\mapsto(1+\epsilon)u$. Write the resulting scalar derivatives as
$P_i,Q_i$. These are formal variations of already selected roots.
They require no resampling of the actual masks.

Suppose the target count $H$ equals one, with abscissa tuple
$\mathbf u^*$. Put $C=\prod_{\mathbf u\ne\mathbf u^*}S(\mathbf u;t)$.
The product rule at this simple zero gives
\[
  N_z=C S_z(\mathbf u^*;t),\qquad
  P_i=C S_{u_i}(\mathbf u^*;t),\qquad Q_i=u_i^*P_i .
\]
Each $P_i$ is nonzero. Fix the other signed points in the unique
completion. Locally the output is translation on the curve followed
by $x$, using $x$ as a local parameter at the signed input.
Translation is an isomorphism and $x$ is unramified at both this
input and $T$ because neither is fixed by negation. The derivative
of the local output with respect to $u_i$ is nonzero. Differentiating
its simple root in $S$ gives the assertion. Thus the original
additive derivatives, with their actual normalization, are
\[
 u_i^*=Q_i/P_i,\qquad N_{s_0}=N_z/S_z(\mathbf u^*;t),
       \qquad N_{s_i}=u_i^* N_{s_0}.
\]
Only one single-tuple form is needed after the split calculation.
The test for this case is $N=0$ and $N_z\ne0$. Upon verifying the
signs and guards it retains the previous $2/9-o(1)$ success bound.

\section{Complete gradient outside the isolated case}
If $N=N_z=0$, the target may have multiple completions at one
abscissa tuple. Its additive gradient can be nonzero. Halve the
choices of each slot, retaining a side with zero scalar norm until
one zero tuple $\mathbf u$ has been located. This uses
$O(m+\sum_i\log n_i)$ scalar calls. Partition the remaining grid
into $m$ boxes by the first coordinate unequal to $u_i$. Multiplying
their scalar norms gives exactly $C=\prod_{\mathbf v\ne\mathbf u}
S(\mathbf v;t)$. Hence $N_{s_0}=C$ and $N_{s_i}=u_i C$, including
when another zero makes $C=0$. The split is valid for the subsets
used internally and for every dyadic subwindow.

A nonzero scalar already answers the relation finder's query. To
evaluate the literal additive derivatives there as well, use the
following two-slot calculation. Fix a prefix $\mathbf w$ of $m-2$
abscissas. Keep the next coordinate $X$ symbolic. For each root
$v$ of the last slot form $F_v(X)=S(\mathbf w,v,t;X)$ and its
derivative with respect to its input $t$. Symmetry identifies it
with the original summation polynomial. Take the subproduct tree of
the polynomials
\[
   F_v(X)+s_0+\sum_{j\le m-2}w_j s_j + Xs_{m-1}+vs_m
\]
with all perturbations, including the target variation, to first
order. Use the fixed-form resultant with
$g(X)=\prod_{u\in U_{m-1}}(X-u)$ and then multiply the scalar
first-order outputs over prefixes. This computes $N$ and every
requested derivative on any input. No pair of choices from the last
two slots is enumerated. The method also handles singular factors.
The $m=1$ case uses the product of linear factors directly.

For rational two-torsion use the paired and individual fixed-point
branches of the preceding note. Each fixed point shifts the target
and contributes its own fixed abscissa times the constant-marker
derivative. The product rule combines all branches. As before the
common local parameter adds $\epsilon$ to each branch's translated
target abscissa. There are at most $4^m$ branches. Alternatively
the asymptotic relation finder may discard seeds containing any
fixed slot point with probability loss at most $3\sum_i\rho_i=o(1)$.

\section{What the hash indices can share}
For a slot of width $\tau=2\Delta$, enumerate the hash preimages
$x_j=(A j+C)\bmod M$ with $M=2^b$. Compare the siblings indexed
by $j$ and $j+\Delta$ for $0\le j<\Delta$. For accepted siblings
the field difference of the abscissas is one of the two constants
$h$, $h-M$ reduced modulo $p$, where $h=A\Delta\bmod M$.
Fix a nonzero curve translation by $Q$. If it realizes a pair for
some choice of signs with a particular difference $s$, then
\[
            K(x_j,x(Q),x_j+s)=0 .
\]
The first term of $K$ supplies $x_j^4$ with coefficient one, and
all other terms have degree at most three.
Each index has a distinct abscissa. Thus one translation can
realize at most eight valid sibling pairs, even with freely chosen
signs and $Q$ selected after seeing the seed. The identity cannot
realize distinct canonical abscissas. Covering $J$ valid pairs this
way takes at least $\lceil J/8\rceil$ translations.

This is relevant to the independent seed law. Let $n_a$ count
two-lift abscissas in $[0,p)$ congruent to $a$ modulo $\tau$. For
uniform odd $A$ and uniform $C$, the first preimage is uniform
modulo $M$. Conditional on it, its sibling is uniform in the
residue class displaced by $\Delta$ modulo $\tau$. Therefore
\[
  \frac{\E J}{\Delta}=\frac{\tau}{M^2}
       \sum_{a\bmod\tau} n_a n_{a+\Delta}.
\]
Uniformly in $a$, $n_a=p/(2\tau)+O(\sqrt p\log p)$. To verify
this, write the class as $a+\tau k$ for a consecutive interval of
$k$. The quadratic character of the squarefree cubic on that
progression has nonzero-frequency complete Fourier sums bounded by
$3\sqrt p$ by the mixed character sum bound. Hasse's bound gives
$2\sqrt p$ at frequency zero. Completing the interval costs
$O(\log p)$ since the sum of the absolute values of its normalized
Fourier coefficients is $O(\log p)$ by the geometric series
formula. The at most three roots of the cubic contribute $O(1)$.
The mixed bound used here is Theorem 5 of
\href{https://pmc.ncbi.nlm.nih.gov/articles/PMC10941333/}
{\emph{Moment subset sums over finite fields}}, with a squarefree
cubic and a nonconstant linear additive phase.
For $\tau\log p=o(\sqrt p)$, including the tuned growing arities,
$\E J/\Delta=(p/(2M))^2(1+o(1))$. Since $p/M\ge1/2$ and
$0\le J\le\Delta$, there is a probability bounded below by an
absolute constant of $J=\Omega(\tau)$. The translation test
concerns reusing the two dyadic halves through fixed elliptic
translations. Other ways of representing the resultant remain open.

\section{End-to-end accounting}
Using fast polynomial products, gcds and extended gcds gives
$\wtO(f+h)$ field operations per resultant and direction once its
two forms are available. Building each individual half form and
its derivatives takes $2^{O(m)}\operatorname{poly}(b,m)$ bit
operations by the quadratic recurrence. Subproduct trees preserve
quasilinear dependence on the sum of formal degrees. Thus the
scalar, its $z$ derivative and all geometric directions take
\[
  2^{O(m)}(D_L+D_R)\operatorname{poly}(b,m)
\]
bit operations. The same bound up to polynomial factors applies
to the full additive gradient at every zero. The two-slot evaluation
on nonzeros has upper bound
\[
  2^{O(m)}\operatorname{poly}(b,m)
    \Big(\prod_{i<m-1}n_i\Big)(n_{m-1}+n_m).
\]
The latter follows because the batched product has degree at most
$2^{m-1}n_m$ in $X$. The two largest slots can be placed last.
Only the former, cheaper bound is needed in the relation finder.

For balanced widths and a balanced split the log of the cheaper
bound is at most $b/2+O(b/m+m+\log b)$, uniformly also on
subwindows. A split chosen after seeing the seeds cannot make
both explicit degrees small. In fact, exactly
\[
       fh=2^{m-1}\prod_i n_i=\tfrac12\prod_i |F_i|
\]
on the paired branch. On prime-order curves of order $r$ the
last product is $(1+o(1))\lambda r$ in probability at the tuned
widths, by the independent slot-size moment estimate in the
previous note. Here $\lambda\in(1/3,2/3]$ and $\log_2 r=b+O(1)$.
Writing the two full coefficient arrays already requires at least
$\sqrt{(1-o(1))\lambda r/2}$ entries with probability tending to
one. The earlier odd-support estimate also prices expanding a
squarefree half polynomial at a fixed balanced split.

There is a retry version for a fresh explicit split per trial. Given
the seed, put $V=\prod_i |F_i|$. At a uniform nonidentity target
the probability of a reachable tuple is at most
$\min(1,V/(r-1))$. Since
$\sqrt V\ge\sqrt{r-1}\min(1,V/(r-1))$, forming the two dense
inputs on each nonempty trial costs $\Omega(\sqrt r)$ per success
after averaging, even with seed-only screening. This argument
assumes that the explicit inputs are formed on admitted trials.
It does not price arbitrary target-dependent screening, implicit
polynomial operations or reuse of an already built table.

The best end-to-end project estimate remains
\[
 \log_2\E T_{\rm solve}\le(26.2933775156+o(1))
              \frac{b}{(\log_2 b)^{1/4}}
\]
under the finite-packet assumption and inherited contraction
hypotheses. The missing saving on this route is a way to evaluate
the final scalar resultant on implicitly represented halves, with
its first-order information or an equivalent isolated-tuple readout,
below the current joint-query cost.

\paragraph{Verification.}
Run \texttt{python -B research/scalar\_norm/verify.py} with
python-flint and NumPy. There are 100 projective resultant tests
with tangents and 567 comparisons of the exact scalar and all
additive derivatives against the previous multiplication matrices.
They include 76 simple targets, four with one repeated zero tuple,
20 with several zero tuples, and 467 nonzero norms (including
empty boxes). Separate checks cover a full rational two-torsion
curve and formal degrees exceeding the characteristic.
On the order-65419 fixture at $(0,1)$, the 2800 abscissa tuples
become 28 left and 100 right tuples, with formal degrees 56 and 400.
The scalar is zero, $N_z=43446$ and $N_{s_0}=16953$ in $\F_{65521}$.
The four coordinate-marker derivatives are 13292, 47523, 12119 and
9163. An independent full-grid product in the verifier checks these
values. On a width-256 hash on the same curve, 29 index siblings
both lift and no curve translation covers more than one of them.
These are finite arithmetic checks. The time comparison above is
asymptotic and uses operation counts.
\end{document}
